\href{}{}% =====================================================================
%  DOSSIER TIPE - Filiere MP/MPI
%  Optimisation du profil d'epaisseur d'un ecran protecteur
%  de smartphone en verre trempe.
%
%  Compilation : pdflatex dossier_tipe.tex   (deux passes pour la table
%  des matieres et les references hyperref).
%
%  Les figures (fig_*.png) sont produites par le script simulation.py.
%  Tant qu'elles n'existent pas, un cadre de substitution est affiche :
%  le document compile donc meme sans avoir lance la simulation.
% =====================================================================
\documentclass[11pt,a4paper]{article}

\usepackage[utf8]{inputenc}
\usepackage[T1]{fontenc}
\usepackage[french]{babel}
\usepackage[margin=2.4cm]{geometry}
\usepackage{amsmath,amssymb}
\usepackage{graphicx}
\usepackage{xcolor}
\usepackage{booktabs}
\usepackage{listings}
\usepackage{caption}
\usepackage{enumitem}
\usepackage[hidelinks]{hyperref}

% ---------------------------------------------------------------------
%  Style des listings Python
% ---------------------------------------------------------------------
\definecolor{codegray}{rgb}{0.5,0.5,0.5}
\definecolor{codegreen}{rgb}{0.0,0.5,0.0}
\definecolor{codeblue}{rgb}{0.1,0.1,0.7}
\definecolor{codeback}{rgb}{0.97,0.97,0.97}

\lstset{
  language=Python,
  backgroundcolor=\color{codeback},
  basicstyle=\ttfamily\scriptsize,
  keywordstyle=\color{codeblue}\bfseries,
  commentstyle=\color{codegreen}\itshape,
  stringstyle=\color{red!60!black},
  numberstyle=\tiny\color{codegray},
  numbers=left,
  numbersep=6pt,
  showstringspaces=false,
  breaklines=true,
  frame=single,
  rulecolor=\color{black!25},
  tabsize=4,
  captionpos=b,
  literate={é}{{\'e}}1 {è}{{\`e}}1 {à}{{\`a}}1 {ê}{{\^e}}1 {û}{{\^u}}1
           {ô}{{\^o}}1 {î}{{\^i}}1 {ç}{{\c c}}1 {É}{{\'E}}1 {-}{-}1
}

% ---------------------------------------------------------------------
%  Encadre "Note jury" : signale les hypotheses simplificatrices fortes
% ---------------------------------------------------------------------
\newcommand{\notejury}[1]{%
  \par\medskip\noindent
  \fcolorbox{red!60!black}{red!4}{%
    \parbox{\dimexpr\linewidth-2\fboxsep-2\fboxrule}{%
      \footnotesize\textbf{\textcolor{red!60!black}{Note jury --- hypothèse à défendre :}}\ #1}}%
  \par\medskip
}

% Inclusion robuste d'une figure produite par la simulation
\newcommand{\figauto}[3]{% #1 fichier  #2 largeur  #3 legende
  \begin{figure}[ht]\centering
  \IfFileExists{#1}{\includegraphics[width=#2\linewidth]{#1}}{%
    \fbox{\parbox[c][3.2cm][c]{#2\linewidth}{\centering\footnotesize
    \textit{[Figure \texttt{#1} : générée en exécutant \texttt{simulation.py}.]}}}}
  \caption{#3}\end{figure}
}

\title{\vspace{-1cm}\textbf{Optimisation du profil d'épaisseur d'un écran\\
protecteur de smartphone en verre trempé}\\[4pt]
\large Minimiser la probabilité de rupture en chute sous contrainte de tactilité}
\author{TIPE --- Filière MP/MPI}
\date{Année 2025--2026}

\begin{document}
\maketitle
\thispagestyle{empty}

\begin{center}
\fbox{\parbox{0.92\linewidth}{\small
\textbf{Thème TIPE associé :} (à adapter au thème de l'année). Le sujet
s'inscrit dans l'étude des \emph{transitions, ruptures et transformations}
d'un système mécanique soumis à une sollicitation brève, et dans la
\emph{recherche d'un optimum sous contraintes}.}}
\end{center}

\tableofcontents
\bigskip
\hrule


% =====================================================================
\section{MCOT --- Mise en cohérence des objectifs du TIPE}
% =====================================================================

\paragraph{Problématique.}
\emph{Quelle répartition spatiale d'épaisseur $e(x,y)$ d'un écran protecteur
en verre trempé minimise la probabilité de rupture lors d'une chute, sous
contrainte de tactilité (épaisseur centrale $e_c \le 0{,}5$~mm) et de masse
raisonnable ?}

\paragraph{Positionnement thématique.}
Le sujet articule trois compétences du programme : (i) la \textbf{mécanique}
d'un choc bref modélisé par un oscillateur amorti et l'approximation
impulsionnelle ; (ii) le traitement \textbf{probabiliste} de la rupture
fragile, naturellement statistique, via une simulation de Monte-Carlo ; (iii)
la recherche d'un \textbf{optimum sous contraintes}. La démarche suit un fil
unique : modèle théorique $\rightarrow$ validation par analogie
électromécanique $\rightarrow$ simulation numérique $\rightarrow$ validation
expérimentale acoustique.

\paragraph{Objectifs (déclinés et vérifiables).}
\begin{enumerate}[leftmargin=2em,itemsep=1pt]
  \item Établir un modèle mécanique du choc reliant l'épaisseur locale $e$ à
        la contrainte de traction maximale $\sigma_{\max}$.
  \item Valider l'ordre de grandeur de ce modèle par une analogie RLC mesurable
        à l'oscilloscope.
  \item Estimer par Monte-Carlo la probabilité de rupture $P_r$ de trois
        familles de profils (uniforme, parabolique, exponentiel).
  \item Optimiser le profil parabolique sous les contraintes de tactilité et de
        masse.
  \item Valider expérimentalement la \emph{carte de rigidité} prédite par une
        analyse acoustique (FFT) de la fréquence propre locale.
\end{enumerate}

\paragraph{Mots-clés.}
\begin{center}\small
\begin{tabular}{ll}
\toprule
\textbf{Français} & \textbf{Anglais} \\
\midrule
Verre trempé & Tempered / strengthened glass \\
Plaque mince (Kirchhoff--Love) & Thin plate (Kirchhoff--Love) \\
Contact de Hertz & Hertzian contact \\
Rupture fragile / loi de Weibull & Brittle fracture / Weibull law \\
Simulation de Monte-Carlo & Monte-Carlo simulation \\
Précontrainte de compression & Compressive prestress \\
\bottomrule
\end{tabular}
\end{center}

% =====================================================================
\section{Introduction et contexte}
% =====================================================================

Les écrans protecteurs (\emph{screen protectors}) en verre trempé dominent
aujourd'hui le marché des accessoires de smartphone, devant les films
plastiques (PET, TPU). Leur intérêt tient à leur dureté de surface élevée
(résistance aux rayures) et à leur toucher proche du verre nu, qui préserve la
sensibilité tactile capacitive. Leur point faible bien connu est le
comportement \emph{aux bords et aux coins} : c'est là que s'amorcent la grande
majorité des fissures observées après une chute, car le coin concentre la
sollicitation et présente une surface de bord souvent moins bien trempée.

\medskip
L'enjeu pratique est double. D'une part, un écran trop fin casse fréquemment ;
d'autre part, un écran uniformément épais dégrade la tactilité et l'esthétique,
et augmente la masse et le coût. L'idée directrice de ce TIPE est qu'\emph{une
épaisseur non uniforme}, plus forte là où la contrainte au choc est élevée
(bords, coins) et plus faible au centre (zone tactile principale), pourrait
réduire fortement la probabilité de rupture à tactilité préservée. Nous
cherchons à quantifier ce gain et à en discuter la faisabilité.

% =====================================================================
\section{Modélisation physique}
% =====================================================================

\subsection{Géométrie et hypothèses}

On modélise l'écran par une plaque mince rectangulaire à coins arrondis :
longueur $L = 150$~mm, largeur $\ell = 70$~mm, rayon de coin $R = 12$~mm, et
épaisseur variable $e(x,y)$. On note $r$ la distance au centre normalisée
($r=0$ au centre, $r=1$ au bord le plus proche).

\medskip\noindent\textbf{Hypothèses du modèle :}
\begin{enumerate}[leftmargin=2em,itemsep=1pt]
  \item \emph{Kirchhoff--Love} (hors-programme, justifiée) : l'épaisseur
        ($\sim 0{,}5$~mm) est très petite devant les dimensions latérales
        ($\sim 100$~mm), donc le cisaillement transverse est négligé et la
        normale reste droite et perpendiculaire au feuillet moyen.
  \item Petites déformations, élasticité linéaire isotrope.
  \item Rupture fragile sans plasticité (le verre ne plastifie pas).
  \item Précontrainte de compression de surface (trempe chimique par échange
        ionique \mbox{K$^+$/Na$^+$}) : la rupture n'apparaît que lorsque la
        traction appliquée dépasse cette précontrainte.
\end{enumerate}

\notejury{Kirchhoff--Love suppose $e \ll L,\ell$ et ignore le cisaillement
transverse. Près des coins très rigides, cette hypothèse est la plus
discutable ; on la conserve car on s'intéresse aux \emph{tendances} comparées
entre profils, pas à la valeur absolue de $\sigma_{\max}$.}

\subsection{Propriétés du matériau}

\begin{center}\small
\begin{tabular}{lll}
\toprule
Grandeur & Valeur retenue & Source / commentaire \\
\midrule
Module d'Young $E$ & $73$~GPa & verre aluminosilicate \cite{corning} \\
Coefficient de Poisson $\nu$ & $0{,}22$ & \cite{corning,timoshenko} \\
Masse volumique $\rho$ & $2450$~kg\,m$^{-3}$ & \cite{corning} \\
Contrainte de compression de surface & $700$--$850$~MPa & trempe chimique \cite{gy} \\
Profondeur de couche (DOL) & $30$--$50\ \mu$m & échange ionique \cite{gy} \\
\bottomrule
\end{tabular}
\end{center}

\subsection{Équation des plaques de Lagrange--Germain}

Pour une plaque d'épaisseur \emph{uniforme}, le déplacement transverse
$w(x,y)$ sous une charge surfacique $q(x,y)$ vérifie l'équation biharmonique de
Lagrange--Germain :
\begin{equation}
D\,\nabla^4 w = q,
\qquad
D = \frac{E\,e^{3}}{12\,(1-\nu^{2})},
\label{eq:biharm}
\end{equation}
où $D$ est la \emph{rigidité de flexion} et $\nabla^4 = \nabla^2\nabla^2$ le
bilaplacien. La dépendance $D \propto e^{3}$ est l'ingrédient central : la
rigidité augmente très vite avec l'épaisseur.

\medskip
Lorsque l'épaisseur varie, $D = D(x,y)$ ne sort plus de l'opérateur et
l'équation devient (forme générale, hors-programme) :
\begin{equation}
\nabla^2\!\big(D\,\nabla^2 w\big)
-(1-\nu)\Big(\partial_{xx}D\,\partial_{yy}w
-2\,\partial_{xy}D\,\partial_{xy}w
+\partial_{yy}D\,\partial_{xx}w\Big) = q.
\label{eq:biharm-var}
\end{equation}
Les termes supplémentaires couplent le gradient d'épaisseur et la courbure.

\notejury{L'équation \eqref{eq:biharm-var} n'est pas résolue analytiquement :
elle justifie seulement que le problème à épaisseur variable est bien posé.
Pour la simulation, on remplace la plaque entière par un \textbf{oscillateur à
1~degré de liberté} équivalent (sous-section suivante), ce qui constitue
l'approximation la plus forte du modèle.}

\subsection{Réduction à un oscillateur amorti à 1 degré de liberté}

On suppose que la réponse au choc est dominée par un mode local de flexion. La
zone qui fléchit réellement est assimilée à une plaque circulaire encastrée de
rayon effectif $a$, chargée en son centre. On obtient un oscillateur
\begin{equation}
m_{\mathrm{eff}}\,\ddot{w} + c\,\dot{w} + k(e)\,w = F(t),
\end{equation}
avec une masse modale $m_{\mathrm{eff}} \propto \rho\,e\,a^{2}$, un
amortissement $c$ (faible pour le verre), et une raideur
\begin{equation}
k(e) = \frac{16\pi D}{a^{2}} = \frac{16\pi}{a^{2}}\cdot\frac{E\,e^{3}}{12(1-\nu^{2})}\ \propto\ e^{3}.
\label{eq:raideur}
\end{equation}
La fréquence propre vaut $\omega_0 = \sqrt{k/m_{\mathrm{eff}}} \propto e/a^{2}$ :
\emph{plus la zone est épaisse, plus elle est rigide et aiguë.} C'est cette
relation que l'expérience acoustique (section~6) cherche à vérifier.

\subsection{Force de choc : approximation impulsionnelle et contact de Hertz}

La force de contact $F(t)$ avec le sol est brève. On la compare à la période
propre. Le \textbf{contact de Hertz} (hors-programme) entre une surface
sphérique de rayon local $R_c$ et le sol donne une durée de contact
\begin{equation}
t_c \simeq 2{,}87\left(\frac{m^{2}}{R_c\,E^{*2}\,v}\right)^{1/5},
\qquad
\frac{1}{E^{*}} = \frac{1-\nu^{2}}{E}+\frac{1-\nu_{\text{sol}}^{2}}{E_{\text{sol}}},
\label{eq:hertz}
\end{equation}
et une force maximale $F_{\max} = k_h\,\delta_{\max}^{3/2}$ avec
$k_h = \tfrac{4}{3}E^{*}\sqrt{R_c}$ et $\delta_{\max}=\big(\tfrac{5mv^2}{4k_h}\big)^{2/5}$.

\medskip
\textbf{Application numérique} (chute de $h=1$~m, $v=\sqrt{2gh}=4{,}4$~m\,s$^{-1}$,
$m=0{,}18$~kg, $R_c=5$~mm) : on obtient $t_c \approx 0{,}14$~ms et
$F_{\max} \approx 21$~kN. La fréquence fondamentale d'une plaque
$150\times70\times0{,}5$~mm vaut $f_{11}\approx 315$~Hz, soit une période
$T_{11}\approx 3{,}2$~ms. Le rapport
\begin{equation}
\frac{t_c}{T_{11}} \approx 0{,}04 \ll 1
\end{equation}
légitime l'\textbf{approximation impulsionnelle} (force assimilée à une
impulsion de Dirac $J\,\delta(t)$) vis-à-vis du mode structural : la plaque
« voit » un transfert quasi instantané de quantité de mouvement
$J = \int F\,dt \approx m\,v\,(1+\varepsilon)$, $\varepsilon$ étant le
coefficient de restitution.

\notejury{L'approximation de Dirac est valable vis-à-vis du \emph{mode global}
($t_c \ll T_{11}$) mais \emph{pas} vis-à-vis des modes très locaux de très
haute fréquence, pour lesquels $t_c$ peut devenir comparable à la période.
La contrainte de contact locale est alors traitée séparément par Hertz.}

\subsection{De la fléche à la contrainte}

Par conservation de l'énergie, une fraction $\eta$ de l'énergie cinétique
$E_k=\tfrac12 m v^{2}$ est stockée en flexion : $\tfrac12 k\,w_0^{2}=\eta E_k$,
d'où la fléche maximale $w_0=\sqrt{2\eta E_k/k}$. Pour une plaque circulaire
encastrée, la contrainte de traction maximale (au bord encastré, en peau) est
\begin{equation}
\boxed{\ \sigma_{\max} = \frac{2E\,e\,w_0}{(1-\nu^{2})\,a^{2}}\,K_c
\ \propto\ \frac{\sqrt{\eta E_k}}{a\,\sqrt{e}}\ }
\label{eq:sigma}
\end{equation}
où $K_c \ge 1$ est un facteur de concentration de contrainte aux coins. Le
résultat clé est $\sigma_{\max}\propto e^{-1/2}$ : \emph{épaissir localement
réduit la contrainte}, ce qui motive l'optimisation.

\notejury{La valeur de la fraction d'énergie $\eta$ ($\approx 0{,}1$) et du
rayon de réponse $a$ sont des paramètres effectifs calés sur l'ordre de
grandeur attendu. Ils n'affectent pas le \emph{classement} des profils, qui est
la conclusion robuste de l'étude.}

\subsection{Critère de rupture fragile}

La rupture est atteinte lorsque la traction de surface dépasse un seuil
$\sigma_c$. Pour un matériau fragile, ce seuil est \emph{statistique} : il
dépend de la population de micro-défauts. On le modélise par une loi de
Weibull, de probabilité de rupture
\begin{equation}
P_r(\sigma) = 1 - \exp\!\left[-\left(\frac{\sigma}{\sigma_\theta}\right)^{m}\right],
\label{eq:weibull}
\end{equation}
avec une contrainte caractéristique $\sigma_\theta \approx 850$~MPa
(intégrant la précontrainte de trempe) et un module de Weibull $m\approx 8$.


% =====================================================================
\section{Analogie électromécanique : validation du modèle}
% =====================================================================

L'oscillateur mécanique amorti est rigoureusement analogue à un circuit
\textbf{RLC série} excité par une brève impulsion de tension. Cette analogie
permet de \emph{mesurer} à l'oscilloscope un régime transitoire dont les
paramètres se transposent au problème mécanique, donc de valider qualitativement
le comportement attendu (régime pseudo-périodique, dépendance de la fréquence à
la « raideur »).

\begin{center}\small
\begin{tabular}{lll}
\toprule
Mécanique & Électrique & Correspondance \\
\midrule
Masse modale $m_{\mathrm{eff}}$ & Inductance $L$ & $m \leftrightarrow L$ \\
Amortissement $c$ & Résistance $R$ & $c \leftrightarrow R$ \\
Raideur $k(e)$ & Inverse de capacité $1/C$ & $k \leftrightarrow 1/C$ \\
Fléche $w$ & Charge $q$ & $w \leftrightarrow q$ \\
Vitesse $\dot w$ & Courant $i$ & $\dot w \leftrightarrow i$ \\
Impulsion $J$ & Échelon de charge initial & $J \leftrightarrow q_0$ \\
\bottomrule
\end{tabular}
\end{center}

L'équation $L\ddot q + R\dot q + q/C = e(t)$ est formellement identique à
celle de l'oscillateur mécanique. La pulsation propre
$\omega_0 = 1/\sqrt{LC}$ joue le rôle de $\sqrt{k/m_{\mathrm{eff}}}$ : faire
varier $C$ (donc « $1/k$ ») revient à faire varier l'épaisseur dans le modèle.

\paragraph{Protocole de mesure.}
\begin{enumerate}[leftmargin=2em,itemsep=1pt]
  \item Monter un circuit RLC série ($L \sim 10$~mH, $C$ variable
        $1$--$100$~nF, $R$ réglable).
  \item Exciter par un créneau bref (GBF) pour simuler l'impulsion de choc.
  \item Visualiser la tension aux bornes de $C$ à l'oscilloscope ; mesurer la
        pseudo-période $T$ et le décrément logarithmique
        $\delta = \tfrac{1}{n}\ln\!\frac{u_k}{u_{k+n}}$.
  \item Faire varier $C$ et vérifier $T \propto \sqrt{C}$, image de
        $T_{\text{méca}} \propto 1/\sqrt{k} \propto e^{-3/2}$.
\end{enumerate}

\paragraph{Comparaison attendue.}
On attend un régime pseudo-périodique amorti, avec une fréquence qui
\emph{augmente} quand on diminue $C$ (analogue d'une épaisseur plus grande,
donc d'une zone plus rigide). Cette tendance confirme qualitativement la
relation $\omega_0 \propto e$ utilisée dans le modèle.

\notejury{L'analogie reproduit la \emph{dynamique} (régime transitoire d'un
système du 2\textsuperscript{nd} ordre) mais pas la \emph{géométrie} de la
plaque ni le critère de rupture. C'est une validation de la brique
« oscillateur », pas du modèle complet.}

% =====================================================================
\section{Simulation numérique}
% =====================================================================

\subsection{Principe et profils testés}

On évalue la probabilité de rupture $P_r$ par une simulation de Monte-Carlo de
$N=5000$ chutes. Pour chaque tirage : un angle de chute, une hauteur, donc une
énergie cinétique ; on en déduit le point d'impact (et donc l'épaisseur locale
selon le profil), la contrainte $\sigma_{\max}$ par \eqref{eq:sigma}, puis on
teste la rupture contre un seuil de Weibull \eqref{eq:weibull} tiré
aléatoirement. Trois profils $e(r)$ sont comparés :
\begin{align}
\text{uniforme :}\quad & e(r)=e_0, \\
\text{parabolique :}\quad & e(r)=e_c+(e_b-e_c)\,r^{2}, \\
\text{exponentiel :}\quad & e(r)=e_c\,e^{\alpha r},\ \ \alpha=\ln(e_b/e_c).
\end{align}

\paragraph{Distribution réaliste de l'angle.}
On n'utilise pas une loi uniforme naïve : les smartphones tombent rarement
parfaitement à plat. On tire l'angle d'impact selon une \textbf{gaussienne
tronquée} centrée à $65^\circ$ (écart-type $15^\circ$, tronquée sur
$[10^\circ,90^\circ]$), ce qui privilégie les impacts proches du bord/coin,
statistiquement les plus dommageables.

\notejury{Le centrage à $65^\circ$ est un choix de modélisation
(comportement « en moyenne sur le bord »). Une étude de sensibilité à ce
paramètre serait une réponse pertinente à une question du jury.}

\subsection{Résultats numériques}

L'exécution du script (annexe~A, $N=5000$, graine fixée) donne les
\emph{ordres de grandeur attendus} suivants :

\begin{center}\small
\begin{tabular}{lcccc}
\toprule
Profil & $e_c$ (mm) & $e_b$ (mm) & $\langle\sigma_{\max}\rangle$ (MPa) & $P_r$ \\
\midrule
Uniforme       & 0{,}50 & 0{,}50 & $\sim 1100$ & $\sim 73\%$ \\
Parabolique    & 0{,}50 & 1{,}20 & $\sim 820$  & $\sim 48\%$ \\
Exponentiel    & 0{,}50 & 1{,}20 & $\sim 790$  & $\sim 44\%$ \\
\midrule
Parabolique \emph{optimisé} & 0{,}50 & 2{,}00 & --- & $\sim 24\%$ \\
\bottomrule
\end{tabular}
\end{center}

\notejury{Ces valeurs sont des \emph{résultats de simulation}, pas des mesures.
La valeur absolue de $P_r$ dépend des paramètres effectifs ($\eta$, $a$,
$\sigma_\theta$) ; c'est la \emph{hiérarchie} entre profils qui constitue le
résultat exploitable.}

On vérifie aussi numériquement la cohérence du choc : $t_c \approx 0{,}14$~ms,
$F_{\max}\approx 21$~kN, $f_{11}\approx 315$~Hz, soit $t_c/T_{11}\approx 0{,}04$,
confortant l'approximation impulsionnelle (section~3.5).

\figauto{fig_carte_pression.png}{0.85}{Carte de contrainte de traction
(profil uniforme, $h=1$~m) : la contrainte est maximale aux coins, minimale au
centre --- ce qui motive d'y placer la matière.}

\figauto{fig_P_de_e.png}{0.7}{Probabilité de rupture en fonction de l'épaisseur
uniforme. La décroissance reflète $\sigma_{\max}\propto e^{-1/2}$ ; la
contrainte de tactilité interdit la zone $e>0{,}5$~mm au centre.}

\figauto{fig_comparaison_profils.png}{0.98}{Profils testés (gauche) et
probabilités de rupture associées (droite) : à $e_c$ identique, épaissir les
bords réduit nettement $P_r$.}

\subsection{Optimisation sous contraintes}

On optimise le profil parabolique sur $(e_c,e_b)$ avec deux contraintes :
\emph{tactilité} ($e_c\le 0{,}5$~mm) et \emph{masse}
($\bar e \le 2{,}5\,\bar e_{\text{unif}}$, $\bar e$ étant l'épaisseur moyenne
pondérée par l'aire). Le balayage donne un optimum à $e_c=0{,}5$~mm,
$e_b=2{,}0$~mm avec $P_r\approx 24\%$ : \emph{les deux contraintes sont
saturées}. L'optimum pousse $e_c$ à son maximum tactile et $e_b$ à la limite de
masse, traduisant un compromis direct « sécurité $\leftrightarrow$ masse ».

\figauto{fig_optimisation.png}{0.85}{Carte d'optimisation $P_r(e_c,e_b)$ du
profil parabolique. Le minimum se situe au coin « $e_c$ maximal, $e_b$ maximal
autorisé par la masse ».}

% =====================================================================
\section{Expérience --- validation acoustique de la rigidité locale}
% =====================================================================

\subsection{Idée et principe}

On ne peut pas casser des dizaines d'écrans pour mesurer $P_r$. On valide
plutôt la brique centrale du modèle : la \textbf{carte de rigidité locale}.
D'après \eqref{eq:raideur}, une zone plus épaisse/plus rigide possède une
fréquence propre plus haute. En excitant la plaque par un impact léger en
différents points et en analysant le son émis par FFT, on reconstruit une carte
de fréquence propre locale $f(x,y)$, image de la rigidité.

\subsection{Protocole détaillé}
\begin{enumerate}[leftmargin=2em,itemsep=1pt]
  \item Poser l'échantillon (écran protecteur réel) sur trois appuis mous
        (mousse) pour s'approcher de conditions libres.
  \item Définir une grille d'au moins $3\times3$ points : centre, milieux des
        bords, coins.
  \item Exciter chaque point par un choc léger et reproductible (petite bille,
        stylet) ; enregistrer le son avec le micro d'un smartphone
        (\texttt{.wav}, $f_s\ge 44{,}1$~kHz).
  \item Pour chaque point, isoler le transitoire après l'attaque, fenêtrer
        (Hann), calculer la FFT et relever la fréquence du pic dominant
        (script, annexe~B).
  \item Construire la carte $f(x,y)$ et la comparer à la carte de contrainte
        théorique.
\end{enumerate}

\subsection{Tableau de mesures (à remplir)}

\begin{center}\small
\begin{tabular}{lcccc}
\toprule
Position & $f$ mesurée (Hz) & $u(f)$ (Hz) & $\sigma_{\text{théo}}$ (MPa) & Rang \\
\midrule
Centre            & \dots & \dots & $\sim 350$ & \dots \\
Milieu bord long  & \dots & \dots & $\sim 600$ & \dots \\
Milieu bord court & \dots & \dots & $\sim 570$ & \dots \\
Coin 1            & \dots & \dots & $\sim 820$ & \dots \\
Coin 2            & \dots & \dots & $\sim 810$ & \dots \\
\bottomrule
\end{tabular}
\end{center}

L'incertitude type sur la fréquence est au moins la résolution spectrale
$u(f)=f_s/N$ ($N$ : nombre d'échantillons de la fenêtre analysée), à compléter
par la dispersion sur plusieurs répétitions.

\subsection{Méthode de comparaison}

On attend une \textbf{corrélation négative} entre la fréquence locale et la
contrainte théorique : une zone \emph{basse fréquence} (souple) correspond à une
zone de \emph{forte contrainte} au choc. Le script calcule le coefficient de
corrélation de Pearson entre les deux cartes ; un $r$ nettement négatif valide
la cohérence rigidité $\leftrightarrow$ contrainte du modèle.

\figauto{fig_carte_acoustique.png}{0.95}{Comparaison carte de fréquence propre
mesurée (gauche) / carte de contrainte théorique (droite). Anti-corrélation
attendue.}

\subsection{Limites honnêtes de l'expérience}

Le son renseigne sur la \emph{rigidité relative} des zones, pas sur le
\emph{seuil de rupture absolu} : deux zones de même rigidité peuvent avoir des
seuils de rupture différents (défauts, qualité de trempe locale). L'expérience
valide donc le \emph{classement} des zones par rigidité, brique nécessaire mais
non suffisante de la prédiction de $P_r$.

% =====================================================================
\section{Discussion critique et limites}
% =====================================================================

\paragraph{Fabricabilité.}
Un gradient \emph{continu} d'épaisseur est difficilement industrialisable : la
trempe chimique par échange ionique se fait en bain et tend à uniformiser la
précontrainte ; un usinage du profil avant trempe est coûteux et fragilise les
bords. Deux alternatives réalistes :
\begin{itemize}[leftmargin=2em,itemsep=1pt]
  \item un profil \textbf{en paliers discrets} (centre fin, bordure plus
        épaisse), techniquement plus accessible qu'un gradient lisse ;
  \item un \textbf{traitement différentiel de la précontrainte} : renforcer
        la compression de surface aux bords plutôt que d'y ajouter de la
        matière (mêmes effets, masse constante).
\end{itemize}

\paragraph{Limites du modèle.}
Le modèle linéaire élastique + oscillateur 1~ddl est une forte réduction d'un
problème de plaque à épaisseur variable et de contact dynamique. La rupture
fragile est par nature statistique (loi de Weibull) et sensible aux défauts de
bord, mal capturés par un modèle déterministe de $\sigma_{\max}$. La valeur
absolue de $P_r$ est donc indicative ; la \emph{hiérarchie} des profils est le
résultat robuste.

\paragraph{Incertitudes à propager.}
Côté acoustique : résolution spectrale, reproductibilité de l'excitation,
amortissement par les appuis. Côté chute : dispersion de l'angle et de la
hauteur, coefficient de restitution, rayon de contact $R_c$. Une analyse de
sensibilité (variation de $\pm$ un paramètre) renforcerait les conclusions.

% =====================================================================
\section{Conclusion}
% =====================================================================

La démarche a suivi un fil cohérent : un \textbf{modèle mécanique} reliant
l'épaisseur locale à la contrainte de traction
($\sigma_{\max}\propto e^{-1/2}$), \textbf{validé dynamiquement} par une
analogie RLC, exploité par une \textbf{simulation de Monte-Carlo}, et dont la
brique « rigidité locale » est \textbf{testée expérimentalement} par analyse
acoustique.

\medskip
\textbf{Réponse à la problématique.} À tactilité préservée ($e_c\le 0{,}5$~mm),
redistribuer la matière vers les bords et les coins réduit fortement la
probabilité de rupture : le profil parabolique optimisé fait passer $P_r$
d'environ $73\%$ (uniforme $0{,}5$~mm) à environ $24\%$, l'optimum saturant les
contraintes de tactilité et de masse. La répartition optimale est donc
\emph{mince au centre, épaisse aux bords}, conformément à la carte de contrainte.

\medskip
\textbf{Ouvertures.} Matériaux à gradient de propriétés (FGM) industriels,
trempe différentielle aux bords, et passage à un véritable critère statistique
de Weibull spatialisé, calibré sur des essais de rupture.

% =====================================================================
\begin{thebibliography}{9}
% =====================================================================
\bibitem{timoshenko} S.~P.~Timoshenko, S.~Woinowsky-Krieger,
\textit{Theory of Plates and Shells}, 2\textsuperscript{e}~éd.,
McGraw-Hill, 1959.

\bibitem{johnson} K.~L.~Johnson, \textit{Contact Mechanics},
Cambridge University Press, 1985.

\bibitem{weibull} W.~Weibull, \textit{A Statistical Distribution Function of
Wide Applicability}, Journal of Applied Mechanics, vol.~18, p.~293--297, 1951.

\bibitem{gy} R.~Gy, \textit{Ion exchange for glass strengthening},
Materials Science and Engineering~B, vol.~149, n\textsuperscript{o}~2,
p.~159--165, 2008.

\bibitem{corning} Corning Incorporated, \textit{Gorilla Glass --- propriétés
mécaniques du verre aluminosilicate} (fiches techniques constructeur).

\bibitem{landau} L.~Landau, E.~Lifchitz, \textit{Théorie de l'élasticité},
Éditions Mir (pour le contact de Hertz et la théorie des plaques).
\end{thebibliography}

% =====================================================================
\appendix
\section{Code Python --- simulation de Monte-Carlo}
% =====================================================================
\lstinputlisting[language=Python]{simulation.py}

% =====================================================================
\section{Code Python --- analyse acoustique par FFT}
% =====================================================================
\lstinputlisting[language=Python]{analyse_fft.py}

\end{document}
